% GENERATED FILE. Edit canonical lessons, then run npm run generate:books.
% canonical-source-hash: fnv1a64:c854e98c29b27dc6
% canonical-lessons: HI-C542-varta, HI-C542-riyayat, HI-C542-sahmati, HI-C542-peskas, HI-C542-madhyasth

\chapter{Coming to Terms}
\label{ch:hi-coming-to-terms}
\hlchaptermodality{542}

\begin{chapteropening}
\textbf{By the end of this chapter:} \emph{I can talk about talks, offers, concessions and go-betweens.}

Complete the last lesson of chapter 542: I can talk about talks, offers, concessions and go-betweens.
\end{chapteropening}

\section[\texorpdfstring{\hi{वार्ता} (vārtā)}{वार्ता (vārtā)}]{\hi{वार्ता} (vārtā) — talks; a negotiation (formal)}

\label{lesson:HI-C542-varta}



\begin{quote}
Before the new one: say the Hindi for a noose, then the Hindi for to tangle.
\end{quote}

\subsection*{You'll want to know: \hi{वार्ता}}
\textbf{\hi{वार्ता}} — \emph{vārtā} — \textquotedblleft{}talks; a negotiation (formal)\textquotedblright{}.

\begin{quote}
\hi{दोनों} \hi{पक्षों} \hi{की} \hi{वार्ता} \hi{अच्छी} \hi{रही।}
\end{quote}

\begin{quote}
\emph{Donõ pakṣoṁ kī vārtā acchī rahī.}
\end{quote}

\begin{quote}
The talks between the two sides went well.
\end{quote}

Feminine.

A formal word, heard in the news.

\begin{cousinweb}
\textbf{\hi{वार्ता}} (\emph{vārtā}) is a tatsama, taken from Sanskrit \emph{vārttā}, \textquotedblleft{}news, an account of what has happened\textquotedblright{} (Monier-Williams); Hindi uses it, in a formal register, for talks between two sides.
\end{cousinweb}

\begin{grammarlens}[title={What you've built}]
One more word for saying what happens if.
\end{grammarlens}

\subsection*{Your turn}
\begin{itemize}
\raggedright
  \item \emph{Say it:} \emph{vārtā}
  \item \emph{Say it:} \emph{vārtā}, once more
  \item \emph{Say it:} \emph{Donõ pakṣoṁ kī vārtā acchī rahī.} — The talks between the two sides went well.
  \item \emph{Recall:} say the Hindi for a noose, then the Hindi for to tangle, then say \emph{vārtā} again
\end{itemize}

\begin{tcolorbox}[breakable,colback=teal!4,colframe=teal!35!black,title={Before you move on}]
What does \hi{वार्ता} mean? (\textquotedblleft{}Talks; a negotiation (formal)\textquotedblright{}.) Say it once more.
\end{tcolorbox}
\section[\texorpdfstring{\hi{रियायत} (riyāyat)}{रियायत (riyāyat)}]{\hi{रियायत} (riyāyat) — a concession; a reduction (in a price)}

\label{lesson:HI-C542-riyayat}



\begin{quote}
Before the new one: say the Hindi for to tangle, then the Hindi for talks.
\end{quote}

\subsection*{You'll want to know: \hi{रियायत}}
\textbf{\hi{रियायत}} — \emph{riyāyat} — \textquotedblleft{}a concession; a reduction (in a price)\textquotedblright{}.

\begin{quote}
\hi{मालिक} \hi{ने} \hi{किराए} \hi{में} \hi{थोड़ी} \hi{रियायत} \hi{दी।}
\end{quote}

\begin{quote}
\emph{Mālik ne kirāe meṁ thoṛī riyāyat dī.}
\end{quote}

\begin{quote}
The landlord took a little off the rent.
\end{quote}

Feminine.

\begin{cousinweb}
\textbf{\hi{रियायत}} (\emph{riyāyat}) is usually traced to Arabic \emph{riʿāyat}, \textquotedblleft{}consideration, concession\textquotedblright{}, through Persian.
\end{cousinweb}

\begin{grammarlens}[title={What you've built}]
One more word for saying what happens if.
\end{grammarlens}

\subsection*{Your turn}
\begin{itemize}
\raggedright
  \item \emph{Say it:} \emph{riyāyat}
  \item \emph{Say it:} \emph{riyāyat}, once more
  \item \emph{Say it:} \emph{Mālik ne kirāe meṁ thoṛī riyāyat dī.} — The landlord took a little off the rent.
  \item \emph{Recall:} say the Hindi for to tangle, then the Hindi for talks, then say \emph{riyāyat} again
\end{itemize}

\begin{tcolorbox}[breakable,colback=teal!4,colframe=teal!35!black,title={Before you move on}]
What does \hi{रियायत} mean? (\textquotedblleft{}A concession; a reduction (in a price)\textquotedblright{}.) Say it once more.
\end{tcolorbox}
\section[\texorpdfstring{\hi{सहमति} (sahmati)}{सहमति (sahmati)}]{\hi{सहमति} (sahmati) — agreement; consent (formal)}

\label{lesson:HI-C542-sahmati}



\begin{quote}
Before the new one: say the Hindi for talks, then the Hindi for a concession.
\end{quote}

\subsection*{You'll want to know: \hi{सहमति}}
\textbf{\hi{सहमति}} — \emph{sahmati} — \textquotedblleft{}agreement; consent (formal)\textquotedblright{}.

\begin{quote}
\hi{दोनों} \hi{दोस्तों} \hi{में} \hi{पूरी} \hi{सहमति} \hi{है।}
\end{quote}

\begin{quote}
\emph{Donõ dostoṁ meṁ pūrī sahmati hai.}
\end{quote}

\begin{quote}
The two friends are in complete agreement.
\end{quote}

Feminine.

\begin{cousinweb}
\textbf{\hi{सहमति}} (\emph{sahmati}) is a modern compound of two Sanskrit words, \emph{saha}, \textquotedblleft{}with\textquotedblright{}, and \emph{mati}, \textquotedblleft{}mind, opinion\textquotedblright{} (Monier-Williams): being of one mind with someone.
\end{cousinweb}

\begin{grammarlens}[title={What you've built}]
One more word for saying what happens if.
\end{grammarlens}

\subsection*{Your turn}
\begin{itemize}
\raggedright
  \item \emph{Say it:} \emph{sahmati}
  \item \emph{Say it:} \emph{sahmati}, once more
  \item \emph{Say it:} \emph{Donõ dostoṁ meṁ pūrī sahmati hai.} — The two friends are in complete agreement.
  \item \emph{Recall:} say the Hindi for talks, then the Hindi for a concession, then say \emph{sahmati} again
\end{itemize}

\begin{tcolorbox}[breakable,colback=teal!4,colframe=teal!35!black,title={Before you move on}]
What does \hi{सहमति} mean? (\textquotedblleft{}Agreement; consent (formal)\textquotedblright{}.) Say it once more.
\end{tcolorbox}
\section[\texorpdfstring{\hi{पेशकश} (peśkaś)}{पेशकश (peśkaś)}]{\hi{पेशकश} (peśkaś) — an offer}

\label{lesson:HI-C542-peskas}



\begin{quote}
Before the new one: say the Hindi for a concession, then the Hindi for agreement.
\end{quote}

\subsection*{You'll want to know: \hi{पेशकश}}
\textbf{\hi{पेशकश}} — \emph{peśkaś} — \textquotedblleft{}an offer\textquotedblright{}.

\begin{quote}
\hi{उसने} \hi{मुझे} \hi{नौकरी} \hi{की} \hi{पेशकश} \hi{की।}
\end{quote}

\begin{quote}
\emph{Usne mujhe naukrī kī peśkaś kī.}
\end{quote}

\begin{quote}
He offered me a job.
\end{quote}

Feminine.

\begin{cousinweb}
\textbf{\hi{पेशकश}} (\emph{peśkaś}) is usually traced to Persian \emph{pēshkash}, \textquotedblleft{}a present, an offering\textquotedblright{}.
\end{cousinweb}

\begin{grammarlens}[title={What you've built}]
One more word for saying what happens if.
\end{grammarlens}

\subsection*{Your turn}
\begin{itemize}
\raggedright
  \item \emph{Say it:} \emph{peśkaś}
  \item \emph{Say it:} \emph{peśkaś}, once more
  \item \emph{Say it:} \emph{Usne mujhe naukrī kī peśkaś kī.} — He offered me a job.
  \item \emph{Recall:} say the Hindi for a concession, then the Hindi for agreement, then say \emph{peśkaś} again
\end{itemize}

\begin{tcolorbox}[breakable,colback=teal!4,colframe=teal!35!black,title={Before you move on}]
What does \hi{पेशकश} mean? (\textquotedblleft{}An offer\textquotedblright{}.) Say it once more.
\end{tcolorbox}
\section[\texorpdfstring{\hi{मध्यस्थ} (madhyasth)}{मध्यस्थ (madhyasth)}]{\hi{मध्यस्थ} (madhyasth) — a mediator; a go-between}

\label{lesson:HI-C542-madhyasth}



\begin{quote}
Before the new one: say the Hindi for agreement, then the Hindi for an offer.
\end{quote}

\subsection*{You'll want to know: \hi{मध्यस्थ}}
\textbf{\hi{मध्यस्थ}} — \emph{madhyasth} — \textquotedblleft{}a mediator; a go-between\textquotedblright{}.

\begin{quote}
\hi{दादा} \hi{जी} \hi{ने} \hi{दोनों} \hi{भाइयों} \hi{के} \hi{बीच} \hi{मध्यस्थ} \hi{का} \hi{काम} \hi{किया।}
\end{quote}

\begin{quote}
\emph{Dādā jī ne donõ bhāiyoṁ ke bīc madhyasth kā kām kiyā.}
\end{quote}

\begin{quote}
Grandfather acted as mediator between the two brothers.
\end{quote}

Masculine or feminine, by the person meant.

\begin{cousinweb}
\textbf{\hi{मध्यस्थ}} (\emph{madhyasth}) is a tatsama, taken whole from Sanskrit \emph{madhyastha}, \textquotedblleft{}being in the middle, between\textquotedblright{} (Monier-Williams).
\end{cousinweb}

\begin{grammarlens}[title={What you've built}]
One more word for saying what happens if.
\end{grammarlens}

\subsection*{Your turn}
\begin{itemize}
\raggedright
  \item \emph{Say it:} \emph{madhyasth}
  \item \emph{Say it:} \emph{madhyasth}, once more
  \item \emph{Say it:} \emph{Dādā jī ne donõ bhāiyoṁ ke bīc madhyasth kā kām kiyā.} — Grandfather acted as mediator between the two brothers.
  \item \emph{Recall:} say the Hindi for agreement, then the Hindi for an offer, then say \emph{madhyasth} again
\end{itemize}

\begin{tcolorbox}[breakable,colback=teal!4,colframe=teal!35!black,title={Before you move on}]
What does \hi{मध्यस्थ} mean? (\textquotedblleft{}A mediator; a go-between\textquotedblright{}.) Say it once more.
\end{tcolorbox}
